% ============================================================================
%  Saussure in HLLSet Algebra --- Part II: Synchronic Linguistics
%  A computational exegesis of Ferdinand de Saussure's
%  Course in General Linguistics, Part Two, in the Boolean lattice
%  and GF(2) ring of HLLSet Algebra
% ----------------------------------------------------------------------------
%  Series:  Saussure in HLLSet Algebra (reverse direction of influence)
%  Build:   pdflatex saussure_hllset_part2_synchronic_linguistics.tex
%           (run twice)
%  Notation follows Part I:
%    saussure_hllset_part1_general_principles.tex
%  and the EWM PUB series.
% ============================================================================
\documentclass[11pt,a4paper]{article}

\usepackage[utf8]{inputenc}
\usepackage[T1]{fontenc}
\usepackage{amsmath,amssymb,amsthm}
\usepackage{amscd}
\usepackage{booktabs}
\usepackage{array}
\usepackage{placeins}
\usepackage{hyperref}
\usepackage[margin=1in]{geometry}
\usepackage{enumitem}
\usepackage{xcolor}
\definecolor{hlblue}{RGB}{0,80,160}
\hypersetup{colorlinks=true,linkcolor=hlblue,urlcolor=hlblue,citecolor=hlblue}

% ---- theorems (aligned with the PUB series and Part I) ----------------------
\newtheorem{principle}{Principle}[section]
\newtheorem{theorem}[principle]{Theorem}
\newtheorem{definition}[principle]{Definition}
\newtheorem{proposition}[principle]{Proposition}
\newtheorem{corollary}[principle]{Corollary}
\newtheorem{remark}[principle]{Remark}
\theoremstyle{remark}
\newtheorem{observation}[principle]{Observation}
\newtheorem{reading}[principle]{Reading}

% ---- shorthands (aligned with Part I and the PUB series) --------------------
\newcommand{\Plane}{\ensuremath{\mathbb{B}}}            % bit plane
\newcommand{\Lat}{\ensuremath{\mathcal{B}}}              % HLLSet (Boolean) lattice
\newcommand{\Tok}{\ensuremath{\mathcal{T}}}              % token space
\newcommand{\HLL}{\ensuremath{H}}                        % an HLLSet
\newcommand{\xor}{\ensuremath{\Delta}}                   % ring addition
\newcommand{\BSS}{\ensuremath{\mathsf{BSS}}}             % BSS measure
\newcommand{\res}{\ensuremath{\mathrm{res}}}             % residual
\newcommand{\spanof}[1]{\ensuremath{\langle #1\rangle}}  % GF(2) span
\newcommand{\supp}{\ensuremath{\mathrm{supp}}}           % support
\newcommand{\Ctx}{\ensuremath{\mathcal{C}}}              % token context
\newcommand{\Ing}{\ensuremath{\mu}}                      % ingestion morphism
\newcommand{\Mat}{\ensuremath{\rho}}                     % materialization morphism
\newcommand{\Lang}{\ensuremath{\mathcal{L}}}             % langue (the lattice system)
\newcommand{\Par}{\ensuremath{\mathcal{P}}}              % parole (usage)
\newcommand{\Sign}{\ensuremath{\mathcal{S}}}             % the sign relation
\newcommand{\Cov}{\ensuremath{\mathrm{cov}}}             % coverage
\newcommand{\Disp}{\ensuremath{\mathcal{D}}}             % the dispatcher
\newcommand{\Sync}{\ensuremath{\mathsf{Syn}}}            % synchronic fiber
\newcommand{\Dia}{\ensuremath{\mathsf{Dia}}}             % diachronic trajectory
\newcommand{\Gen}{\ensuremath{\gamma}}                   % basis generation
\newcommand{\cover}{\ensuremath{\mathrm{cover}}}         % frame cover
\newcommand{\Val}{\ensuremath{\mathcal{V}}}              % value (differential)
\newcommand{\Synt}{\ensuremath{\mathsf{Synt}}}           % syntagmatic reading
\newcommand{\Ass}{\ensuremath{\mathsf{Ass}}}             % associative reading

% ---- title -----------------------------------------------------------------
\title{\textbf{Saussure in HLLSet Algebra}\\[4pt]
\large Part II: Synchronic Linguistics\\[4pt]
\large A Computational Exegesis of the \textit{Course in General Linguistics},\\
Part Two, in the Boolean Lattice and GF(2) Ring of HLLSet Algebra\\[6pt]
\normalsize Second paper in the series \textit{Saussure in HLLSet Algebra}}

\author{Alex Mylnikov \and Alexandr Mylnikov Jr. \and Deependra Kumar}
\date{\today}

\begin{document}
\maketitle

\begin{abstract}
Part II of the series reads Saussure's Part Two, \emph{Synchronic
Linguistics}, out of HLLSet Algebra. The center of that part of the
\textit{Course} is the differential theory of value, and the center of the
exegesis is the hard key: value lives in the relations between nodes, and
relative to a basis frame
$F=(B_1,\dots,B_k)$ the relational position of every node of the span has a
unique coordinate vector
$c\in\mathrm{GF}(2)^k$---the exact compression of its \emph{value}. We prove
that
value is a complete invariant on the span---within one synchronic fiber,
identity and value coincide, and the difference of any two signs is itself a
sign,
$X \xor Y = \bigoplus_{i:c_i \neq d_i} B_i$---so Saussure's ``in the language
itself there are only differences, without positive terms'' holds literally:
the only positive terms in the description are the basis vectors, and these
themselves are arbitrary by the coordinate-freedom theorem of Part~I.
The earlier chapters of Part Two are re-founded in order: concrete entities
are sign pairs $(\Ing,\Mat)$ delimited by the coverage round-trip; identity is
content addressability (the $8{:}45$ Geneva--Paris express is one node under
many material compositions); reality is position in the frame. The two kinds
of relation are then separated: syntagmatic relations (in praesentia) are the
document coverage reading along the materialized chain, and associative
relations (in absentia) are the BSS profile against the frame. The mechanism
of language
becomes the dispatcher composing the two readings; the subdivisions of grammar
(morphology, syntax, lexicology) become one algebra over one plane, which is
Saussure's claim that they interpenetrate, proved structurally. The paper
closes with the complete dictionary of Part Two and the residual that Part~III
will take up: diachronic change.
\end{abstract}

\noindent\textbf{Keywords:} Saussure; synchronic linguistics; linguistic
value; differences without positive terms; concrete entity; identity;
reality; syntagmatic relations; associative relations; hard key; soft key;
BSS; basis frame; GF(2) ring; Boolean lattice; HLLSet; computational
exegesis.

% ============================================================================
\section{Introduction}
\label{sec:intro}
% ============================================================================

\subsection{Where Part II stands in the series}

Part I of this series \cite{part1} interpreted Part One of the
\textit{Course}---the sign, arbitrariness, linearity, immutability and
mutability, and the static/evolutionary distinction---as theorems and
invariants of HLLSet Algebra. It deliberately stopped at the door of
value: the signified was assigned to the \emph{differential value} of a node,
but value itself was left undefined. Part Two of the \textit{Course},
\emph{Synchronic Linguistics}, is where Saussure defines it, and Part II of
this series is where the algebra executes it.

The thesis of this paper is one sentence: \textbf{linguistic value lives in
the relations between nodes, and its exact compression is the hard key}.
Relative to a basis frame $F$, a node $X \in \spanof{F}$ is a unique GF(2)
sum of basis vectors, and its coordinate vector $c$ is the compressed form of
what Saussure calls its value. Everything else in Part Two---concrete
entities, identities, realities, syntagmatic and associative relations, the
mechanism of language, the subdivisions of grammar---is re-founded around
this single definition.

The Introduction of the \textit{Course}---langue and parole, and the object
of linguistics---was read in the EWM direction in \cite{langue2026}; Part I
promised its reverse reading here. It is delivered in
Section~\ref{sec:mechanism}: langue is the shared lattice read through the
current frame, parole is the stream of host outputs that actualizes it, and
the dispatcher is the executive that chooses which parole speaks. The
distinction is re-founded from the algebra side, not assumed.

\subsection{Method, unchanged}

The method is that of Part I \cite{part1}: no metaphor, principles as theorems
or invariants, and ambiguity recorded as data. The translation note of Part I
applies unchanged: we cite Roy Harris's translation \cite{harris1983}, keep
the standard pair \emph{signifier}/\emph{signified} (Harris:
\emph{signal}/\emph{signification}), and retain the Russian edition
\cite{saussure1977} for verification. References to the \textit{Course} are by
part, chapter, and section.

\subsection{Roadmap}

Section~\ref{sec:prelim} fixes the standing definitions inherited from Part I.
Section~\ref{sec:entities} treats chapters I--III of Part Two (concrete
entities, identities, realities). Section~\ref{sec:value} is the center:
chapter IV, linguistic value, re-founded as the hard key and the two
``differences'' theorems. Section~\ref{sec:relations} separates syntagmatic
and associative relations (chapter V). Section~\ref{sec:mechanism} covers
chapters VI--VIII (mechanism, grammar, abstract entities). Section
\ref{sec:dictionary} collects the dictionary; Section~\ref{sec:adds} states
what Part II adds to Part I; Section~\ref{sec:conclusion} concludes and hands
the diachronic residual to Part III.

% ============================================================================
\section{Preliminaries: standing definitions from Part I}
\label{sec:prelim}
% ============================================================================

We restate, with their labels, the definitions Part I established; proofs and
motivation are in \cite{part1}.

\begin{definition}[Plane, lattice, ring]
\label{def:plane}
The \textbf{plane} is $\Plane = \{0,\dots,32767\}$; the \textbf{HLLSet
lattice} is $\Lat = \mathcal{P}(\Plane)$, a Boolean lattice under
$\cup,\cap,\subseteq$ and a Boolean ring under $\xor$ (written $\Delta$) and
$\cap$. Cardinalities are $|X|$; the complement of $X$ in $\Plane$ is
$\overline{X}$.
\end{definition}

\begin{definition}[IICA morphism; ingestion; materialization]
\label{def:iica}
$\Ing : \Tok \to \Lat$ maps each token to its \textbf{bit-set} (the bit
positions it activates, a node of the lattice) and is an \textbf{IICA
morphism} if it is idempotent (deterministic), immutable, and
content-addressed; the bit positions are HLLSet-agnostic, and the morphism
extends to finite token sets by
$\Ing(X)=\bigcup_{t\in X}\Ing(\{t\})$. \textbf{Materialization}
$\Mat : \Lat \dashrightarrow \Ctx$ is the partial inverse that restores an
ordered token context through the door, by reading each set bit and
disambiguating through the LUT fiber of that bit which tokens set it.
(The LUT, multi-hashing, and the global gates $G_1,G_2,G_3$ are introduced
in Part I, Section~2.)
\end{definition}

\begin{definition}[BSS; document coverage; frames; generations; keys]
\label{def:frame}
Against a reference $B_w \subseteq \Plane$, the \textbf{directed soft key} is
$\BSS_w(X) = |X \cap B_w| / |B_w|$. For a context $X$ read against a document
$D$, the \textbf{document coverage reading} is
\[
  N_D = |D \setminus X|, \qquad R_D = |D \cap X|, \qquad \Cov = R_D/|D|.
\]
A \textbf{basis frame} is a reduced (RREF) GF(2) basis $F=(B_1,\dots,B_k)$ of
a span of observed nodes, stamped with a \textbf{generation} $\Gen$ that
increments on recomputation. Relative to $F$, each $X \in \spanof{F}$ has a
unique \textbf{hard key} $c \in \mathrm{GF}(2)^k$ with
\[
  X \;=\; \bigoplus_{i : c_i = 1} B_i,
\]
the \textbf{soft key} $w_i = \BSS_{B_i}(X)$, and the \textbf{spill}
$s = |X \setminus \cover(F)|$.
\end{definition}

\begin{definition}[Synchronic fiber; diachronic trajectory]
\label{def:syn-dia}
The \textbf{synchronic fiber} at time $t$ is $\Sync(t) = (\Lat, F(t),
\Gen(t))$; the \textbf{diachronic structure} is the family $\Dia = \{\Sync(t)
: t \in T\}$ with the recorded transitions between consecutive fibers.
\end{definition}

\begin{definition}[The sign]
\label{def:sign}
The \textbf{sign} is the pair $\Sign = (\Ing,\Mat)$; for a node $A \in \Lat$,
the signifier is the bit-set $A$ and the signified is its differential value
plus the materialized context $\Mat(A)$. The round-trip is guarded:
$\Mat(\Ing(\Ctx)) = \Ctx$ up to order iff $\Cov(\Ctx; D)=1$, where the
coverage is read against the document $D$ from which the context $\Ctx$ was
drawn \cite{langue2026}.
\end{definition}

Throughout, $F=(B_1,\dots,B_k)$ is a fixed frame of a fixed fiber
$\Sync(t)$ unless stated otherwise; ``value'' is defined in
Section~\ref{sec:value} and the earlier sections use it provisionally.

% ============================================================================
\section{Concrete entities, identities, realities}
\label{sec:entities}
% ============================================================================

\subsection{Concrete entities}
\label{sec:concrete}

Saussure rejects the idea that the units of a language are given in advance,
either as sounds or as ideas. The linguistic unit is a \emph{concrete entity}:
a segment of the speech chain that unites a concept and a sound pattern, and
the hard problem of chapter II is \emph{delimitation}: the chain is
continuous, and units must be cut out of it \cite[Part II, ch.~II]{harris1983}.

\begin{reading}[Concrete entity]
\label{read:concrete}
A concrete entity is a sign pair: a node $A \in \Lat$ together with its
materialized context $\Ctx_A = \Mat(A)$. The speech chain is the ordered token
stream; ingestion compiles a segment of it to a node; delimitation is the
round-trip law---a segment is a well-delimited unit exactly when
materialization restores it, that is, when coverage is $1$:
\[
  \Mat(\Ing(\Ctx)) = \Ctx \quad \text{up to order}
  \qquad\Longleftrightarrow\qquad
  \Cov(\Ctx) = 1 .
\]
\end{reading}

\begin{remark}[Delimitation as coverage]
Saussure's delimitation problem becomes a measured decision. A candidate
segment with coverage $<1$ is not yet a unit of the system: part of it spills
outside the frame, and the spill $s$ says how much. The algebra does not
pretend delimitation is easy; it says exactly how hard it is, segment by
segment.
\end{remark}

\subsection{Identities}
\label{sec:identities}

Saussure's example is the $8{:}45$ Geneva--Paris express: the train is ``the
same'' each day although its engine, carriages, and personnel differ. Identity
is not material; it is position in a system. The same point applies to a
street demolished and rebuilt, and to a chess piece lost and replaced by
another \cite[Part II, ch.~III]{harris1983}.

\begin{reading}[Identity]
\label{read:identity}
Identity is content addressability---a property of the morphism, not of any
particular implementation such as a hash. Two material occurrences are the
same sign exactly when they ingest to the same node:
\[
  \text{occurrence}_1 \sim \text{occurrence}_2
  \qquad\Longleftrightarrow\qquad
  \Ing(\text{occurrence}_1) = \Ing(\text{occurrence}_2).
\]
The $8{:}45$ express is one node $A$ under many material compositions: the
tokens (engine, carriages, personnel) differ, the node does not. The
express's identity persists because ingestion is deterministic and the plane
is shared.
\end{reading}

\begin{proposition}[Identity is content addressability, and it is stable]
\label{prop:identity}
For any two token contexts $X, Y \subseteq \Tok$,
\[
  X \sim Y
  \qquad\Longleftrightarrow\qquad
  \Ing(X) = \Ing(Y)
\]
where $\sim$ is sign identity: two occurrences are the same sign exactly
when they ingest to the same node. Because $\Ing$ is deterministic (same
content, same node), immutable (the node is never rewritten), and
content-addressed (the node is named by its content alone), the
equivalence is stable across hosts, tokenizers, and times. Its
\emph{value} (Section~\ref{sec:value}) may change with the frame; its
identity cannot.
\end{proposition}

\begin{remark}[Content addressability is not hashing]
\label{rem:not-hashing}
Readers of the EWM papers will associate content addressability with the
reference hash function. The hash is a convenience, not the concept. Any
algorithm or set of rules that (i) recovers the address from the content and
(ii) keeps different contents at different addresses can serve as the IICA
morphism. Because the HLLSet lattice is a Boolean lattice---a Boolean
algebra---the most general such rule is a \emph{logical} one: choose a family
of predicates
\[
  P_1,\dots,P_m : \Tok \to \{\top,\bot\},
\]
and set
\[
  \Ing(t) \;=\; \{\, i : P_i(t) = \top \,\}.
\]
The address of a content is then the truth table of the content under the
chosen rules: the set of bit positions at which the content satisfies the
corresponding predicate. Distinct contents occupy distinct nodes exactly when
the predicate family separates them. Any Boolean expression over the plane
that is associated with a specific content, and that keeps different contents
at different addresses, is therefore a legitimate ingestion morphism. The
hash is one convenient such rule set, preferred in the reference
implementation only for its even distribution of values over the plane---not
for any linguistic reason.
\end{remark}

\subsection{Realities}
\label{sec:realities}

Saussure's chapter III culminates in the claim that the \emph{reality} of a
linguistic unit is not its substance but its place in the system: what is
real in a language is what speakers recognize as a value among other values
\cite[Part II, ch.~III]{harris1983}.

\begin{reading}[Reality]
\label{read:reality}
The reality of a node in a fiber is its position in the frame. Recognition
has two modes. The \emph{hard key} $c$ recognizes a node by matching a
linear composition over the ring basis: it exists exactly when
$X \in \spanof{F}$, and it is the node's reality in full. The \emph{soft
key} profile $w$ recognizes by matching bit-vectors directly---each
coordinate is the BSS overlap of $X$ with one basis vector---so it is
always defined, for every node; when $X \notin \spanof{F}$ the soft key is
the partial reality, and the spill $s$ is the size of the part that has no
reality in this system yet.
\end{reading}

\begin{observation}[Substance is not reality]
A bit position, a token, and a host codebook are substances; none of them is a
linguistic reality. Reality begins at the first coordinate: the hard key $c$
is what the system itself can say about the node, and nothing less is a unit
of langue. This is the algebra's version of Saussure's anti-substantialism.
\end{observation}

% ============================================================================
\section{Linguistic value}
\label{sec:value}
% ============================================================================

\subsection{Thought and sound coupled}
\label{sec:thought-sound}

Saussure opens chapter IV with the image of the sheet of paper: thought is the
front, sound the back; one cannot cut the front without cutting the back
\cite[Part II, ch.~IV, \S1]{harris1983}. Language is the domain where the two
orders of difference---conceptual and material---are articulated together.

\begin{reading}[The sheet of paper]
\label{read:paper}
The sheet is the plane $\Plane$. The sound side of a sign is its bit-set
$A \subseteq \Plane$; the thought side is its value in the frame (its hard
key $c$). The coupling of the two faces rests on an invariant substrate: the
assignment of bit positions to contents. A bit position assigned to a
token---more exactly, to the collection of tokens that set it---stays the
same whatever HLLSet we form; it is \emph{HLLSet-agnostic}. What the lattice
changes is never a position but a node: a set of positions. A join
$A \cup B = C$ is a new HLLSet into which the tokens and their bit positions
pass together---which is to say, they stay exactly where they were. The two
faces are two readings of one node, not two objects glued together; and the
sheet itself is never cut.
\end{reading}

\begin{remark}[Bit positions do not move; nodes do]
\label{rem:bit-agnostic}
One point deserves plain statement, because the set-theoretic notation can
obscure it. When the lattice forms $A \cup B = C$, the new node $C$ is
simply the set of positions that already belonged to $A$ or to $B$. No
position migrates and no token changes its address: the bits that enter
$C$ were already where they are. The sound side of the sign is therefore
coupled to the thought side at the level of the \emph{node}---a node and
its value are formed together---while the underlying sheet, the position
assignment fixed once by ingestion, is common to all nodes and never
changes. Saussure's ``one cannot cut the front without cutting the back''
becomes: one cannot form a new node without forming, at the same time, a
new value; but one never cuts the positions themselves.
\end{remark}

\subsection{Conceptual value}
\label{sec:conceptual-value}

Saussure's first move is negative: the conceptual side of a sign is not a
pre-existing concept but a value arising from the presence of other signs in
the same system. The stock example is French \emph{mouton} covering both
English \emph{sheep} and \emph{mutton}: the same concept has different values
because the neighbouring terms differ \cite[Part II, ch.~IV, \S2]{harris1983}.
The example is relational through and through: \emph{mouton}'s value is not in
the animal but in the fact that the French system lacks the opposition English
draws. Value lives in relations between nodes, not in nodes.

\begin{reading}[Conceptual value is relational]
\label{read:conceptual}
The conceptual value of a node is not in the node; it is in the relations
between nodes. The hard key $c \in \mathrm{GF}(2)^k$ is not the value itself
but its coordinate compression in the current frame: an exact and economical
projection of the node's relational position against the oppositions
$B_1,\dots,B_k$. The value proper is the node's full relational neighbourhood
in the lattice---which nodes it contains and is contained in, which it joins,
meets, and differs from, and which co-occur with it in materialized contexts.
The same node has different values in different frames, exactly as the same
animal has different values in the French and English lexicons; what the
frames change is the relational reading, not the node.
\end{reading}

\begin{definition}[Relational structure; co-occurrence profile]
\label{def:relational}
A \textbf{relational structure} in the lattice is a pair $G = (N, R)$, where
$N \subseteq \Lat$ is a finite set of nodes and $R$ is a set of binary
relations on them---containment, disjointness, xor-adjacency, agreement of
frame coordinates, and the like. The \textbf{co-occurrence profile} of a node
$n \in N$ is the vector of co-occurrence counts of $n$ with the other nodes
of $N$ across materialized contexts (the $R_D$ values of the document
coverage reading, accumulated over windows).
\end{definition}

\begin{observation}[Similar concepts: two maps, not one]
\label{obs:similar-concepts}
Two subgraphs $G_1 = (N_1, R_1)$ and $G_2 = (N_2, R_2)$ of the lattice may
represent similar concepts only if two bijections exist,
\[
  \varphi : N_1 \leftrightarrow N_2,
  \qquad
  \psi : R_1 \leftrightarrow R_2,
\]
which preserve the relational structure,
\[
  (n, n') \in r
  \quad\Longleftrightarrow\quad
  (\varphi(n), \varphi(n')) \in \psi(r)
  \qquad
  \text{for all } n, n' \in N_1,\; r \in R_1,
\]
and under which the co-occurrence profiles match:
$\mathrm{co}(\varphi(n)) \approx \mathrm{co}(n)$ for every $n \in N_1$. The
word ``may'' is carried by the second condition: relational shape alone does
not decide, because two subgraphs can be structurally identical yet co-occur
differently in their respective contexts. Similarity of concepts is therefore
similarity of structure \emph{together with} similarity of co-occurrence---a
graph isomorphism plus a matching of the statistics of usage.
\end{observation}

\begin{remark}[``Ah, they are conceptually the same!'']
\label{rem:aha}
The everyday exclamation that accompanies an analogy is the datum this
section formalizes. ``Ah, they are conceptually the same!'' says three things
at once: (i) the concept is not a property of the things involved but of the
structure of their relations; (ii) the sameness is therefore a structural
sameness---a matching of relations; and (iii) the matching is not given for
free: it requires carrying the elements of one event onto the elements of the
other, element by element and relation by relation, which is exactly the pair
of bijections of Observation~\ref{obs:similar-concepts}. Conceptual identity
is not the identity of objects but the existence of a structure-preserving
mapping between two events, together with a matching of their co-occurrence.
The concept, in this sense, \emph{is} the mapping: the act of recognition
that one event can be read as another. Saussure's later treatment of analogy
belongs to the diachronic half of the \textit{Course}; what Part~III will
show is that the same pair of bijections, applied across generations, is also
the engine of change.

\end{remark}

\begin{theorem}[Value is a complete invariant on the span]
\label{thm:complete}
Let $F$ be a frame and $X,Y \in \spanof{F}$ with hard keys $c,d$. Then
\[
  X = Y \qquad\Longleftrightarrow\qquad c = d,
\]
and the difference of any two signs is itself a sign:
\[
  X \xor Y \;=\; \bigoplus_{i : c_i \neq d_i} B_i .
\]
Within one fiber, identity and value coincide: two signs are distinct exactly
when their values are distinct.
\end{theorem}

\begin{proof}
Both directions of the equivalence are the uniqueness of RREF coordinates over
GF(2) (Definition~\ref{def:frame}). For the difference, expand both nodes in
the basis and use $B_i \xor B_i = \varnothing$, so only coordinates that
differ survive.
\end{proof}

\begin{remark}[The basis is an ontology; value is itself an HLLSet]
\label{rem:ontology}
Theorem~\ref{thm:complete} shows that the basis is an \emph{ontology} in
exactly the sense in which a system of labels is: a finite set of relations
among nodes, not a claim about semantic correctness. Semantic correctness is
a judgment for speakers; the lattice records only relations. A
human-readable ontology (meaningful labels) and the ring basis are therefore
isomorphic presentations of the same object---both are systems of
oppositions, the labels being one presentation and the basis vectors the
other. And because any collection of HLLSets has a natural presentation as a
single HLLSet (its union), value itself---conceptual as well as material---is
an HLLSet, stamped with its provenance.
\end{remark}

\begin{corollary}[The value table is the lexicon]
\label{cor:lexicon}
The lexicon of a fiber is the value table: the list of nodes of the span with
their coordinate vectors. No two rows repeat a vector; the vector is the
word's entire conceptual content \emph{within that frame}---its relational
position, compressed without loss on the span---and the table is the system
of values made literal.
\end{corollary}

\begin{table}[htbp]
\centering
\renewcommand{\arraystretch}{1.2}
\begin{tabular}{@{}llc@{}}
\toprule
Node & Expression in the frame & Hard key (value) \\
\midrule
$\varnothing$ & the zero node & $(0,0,0)$ \\
$X$ & $B_1 \xor B_3$ & $(1,0,1)$ \\
$Y$ & $B_2$ & $(0,1,0)$ \\
$Z$ & $B_1 \xor B_2 \xor B_3$ & $(1,1,1)$ \\
$X \xor Y$ & $= Z$ & $(1,1,1)$ \\
\bottomrule
\end{tabular}
\caption{A miniature lexicon: three oppositions, four nodes, four values.
Identity is value; difference is ring addition.}
\label{tab:lexicon}
\end{table}

\subsection{Material value}
\label{sec:material-value}

Saussure then applies the same negative move to the signifier: what matters in
a sound is not its positive quality but the differences that distinguish it
from all other sounds \cite[Part II, ch.~IV, \S3]{harris1983}. The signifier
is differential on its own side.

\begin{reading}[Material value]
\label{read:material}
The material value of a node is not which bits are set---bits are anonymous
substance---but the difference pattern between it and the other nodes. The
exact material difference of two nodes is their symmetric difference
$X \xor Y$, a node of the same lattice; the relaxed difference is the soft
key profile. A bit position that distinguishes nothing carries no material
value, however ``positive'' it looks.
\end{reading}

\begin{proposition}[The material difference is the ring sum]
\label{prop:material-difference}
For all $X,Y \in \Lat$, the exact difference between the two signifiers is
\[
  X \xor Y \;=\; (X \setminus Y) \cup (Y \setminus X),
\]
and $X \xor Y = \varnothing$ iff $X = Y$. The material difference is thus a
node of the same lattice, and the lattice is closed under differences: the
system can \emph{say} the difference between any two signs, in its own
language.
\end{proposition}

\begin{remark}[Difference is not absolute: it is difference in a given context]
\label{rem:difference-relative}
The ring sum $X \xor Y$ is an absolute node of the lattice, but \emph{what the
difference consists in} is not absolute. It is read off in a given context,
exactly as coordinates in a Boolean ring are read off in a given RREF basis.
In a frame $F$, the difference of two signs is their coordinate difference
\[
  c \xor d \;\in\; \mathrm{GF}(2)^k,
\]
the set of oppositions on which the two values disagree, and only through it
does the absolute node acquire a reading:
\[
  X \xor Y \;=\; \bigoplus_{i : c_i \neq d_i} B_i .
\]
The same pair of nodes, in another frame, differs in other oppositions; and
outside the span the difference has no exact coordinates at all, only the
relaxed soft profile. Difference, like value, is a function of
$(\text{node}, \text{frame})$, never of the node alone. Saussure's differences
are therefore \emph{contextual} differences, and the algebra makes the context
explicit: it is the current frame.
\end{remark}

\subsection{The sign in its totality: only differences}
\label{sec:totality}

Saussure's conclusion for the whole sign is the most quoted sentence of the
\textit{Course}: ``in the language itself there are only differences. Even
more important: a difference generally implies positive terms between which
the difference is set up; but in language there are only differences
\emph{without positive terms}'' \cite[Part II, ch.~IV, \S4]{harris1983}.

\begin{theorem}[Differences without positive terms]
\label{thm:differences}
Let $F$ be a frame. Every node $X \in \spanof{F}$ is a sum of basis vectors,
\[
  X \;=\; \bigoplus_{i : c_i = 1} B_i,
\]
and nothing else is said about it: its entire description is the difference
pattern $c$. The basis vectors themselves are not positive terms either, since
any basis of the same span describes the same system of values up to a
linear isomorphism of $\mathrm{GF}(2)^k$, and no structural property
prefers one basis over another (coordinate freedom, Part I
\cite{part1}). Hence, within the span, there are only differences: each sign
is a position in a field of oppositions, and the field has no absolute origin.
\end{theorem}

\begin{remark}[The zero node]
The node $\varnothing$ is the neutral element of the ring: the absence of all
oppositions. It is not a positive term---it is what every difference is
measured from, and it says nothing by itself. Saussure's ``without positive
terms'' finds its exact counterpart: the zero of the GF(2) vector space is
not a term of the language but the point from which all values are
coordinates.
\end{remark}

\subsection{Value measured: the axiometer}
\label{sec:axiometer}

Part I left value as a definition; the algebra also provides the instruments
that \emph{measure} it.

\begin{proposition}[Hard key, soft key, and spill partition the readings]
\label{prop:measure}
For a node $X$ and frame $F$:
\begin{enumerate}[leftmargin=*,nosep]
  \item if $X \in \spanof{F}$, the \textbf{hard key} $c$ is the exact value;
  \item in all cases, the \textbf{soft key} $w_i = \BSS_{B_i}(X)$ is the
    relaxed value profile, one coordinate per opposition;
  \item the \textbf{spill} $s = |X \setminus \cover(F)|$ is the size of the
    part of $X$ that has no value yet.
\end{enumerate}
The three together are the axiometer of the system: they answer, for every
node, the question ``what is this sign worth, and how much of it remains
unvalued?''
\end{proposition}

\begin{remark}[The residual of measurement]
The spill is not a defect; it is the boundary of the system of values. When
the spill accumulates, the frame pivots (generation $\Gen+1$) and values are
recomputed---which is the diachronic event whose systematic study belongs to
Part III. Saussure's synchronic Part Two stops exactly where the spill is
zero; the algebra can say where that boundary lies.
\end{remark}

% ============================================================================
\section{Syntagmatic and associative relations}
\label{sec:relations}
% ============================================================================

\subsection{The two relations}
\label{sec:two-relations}

Chapter V of Part Two distinguishes the two kinds of relation that organize a
language. \emph{Syntagmatic} relations hold \emph{in praesentia}: between
terms that are combined in the chain. \emph{Associative} relations hold
\emph{in absentia}: between a term and the others that could replace it, in
memory \cite[Part II, ch.~V, \S1]{harris1983}.

\begin{reading}[The two relations]
\label{read:relations}
The syntagmatic reading of a state is along the materialized chain; the
associative reading is against the frame.
\end{reading}

\begin{definition}[Syntagmatic and associative readings]
\label{def:readings}
Fix a fiber $\Sync(t)$ with frame $F$. For a context $X$ read against a
document $D$, the \textbf{syntagmatic reading} is
\[
  \Synt(X; D) \;=\; (N_D, R_D, \Cov)
  \;=\;
  \bigl(|D \setminus X|,\; |D \cap X|,\; |D \cap X|/|D|\bigr).
\]
For the same $X$ against the frame, the \textbf{associative reading} is the
profile
\[
  \Ass(X; F) \;=\; \bigl(\BSS_{B_1}(X), \dots, \BSS_{B_k}(X)\bigr),
\]
replaced by the hard key $c$ when $X \in \spanof{F}$.
\end{definition}

\subsection{Syntagmatic relations}
\label{sec:syntagmatic}

\begin{proposition}[The syntagmatic axis is the chain, restored]
\label{prop:syntagmatic}
The syntagmatic axis is the axis of combination. In the algebra it has three
layers: (i) the \emph{chain} is the ordered token context restored by
materialization $\Mat$; (ii) \emph{composition} is the join: the syntagm
$X \cup Y$ is the union of the constituent nodes; (iii) \emph{contrast in
the chain} is the ring sum $X \xor Y$, the exact phrase-level difference
between two syntagms. The coverage numbers $N_D, R_D, \Cov$ are the
syntagmatic measurement: novelty $N_D$ is what the document holds beyond the
syntagm (the part the chain does not cover), retention $R_D$ is what the
syntagm confirms, coverage is how much of the document the syntagm accounts
for.
\end{proposition}

\begin{observation}[Solidarities are coverage guards]
Saussure's syntagmatic solidarities---the mutual determination of elements in
a combination---are the coverage guard of the round-trip law. A syntagm holds
together exactly when its parts delimit one another, that is, when
materialization of the joined node restores the chain. Below coverage $1$,
the syntagm has spill: it is not yet a unit of the system.
\end{observation}

\subsection{Associative relations}
\label{sec:associative}

\begin{proposition}[The associative axis is the frame profile]
\label{prop:associative}
The associative axis is the axis of substitution. In the algebra it is the
profile $\Ass(X; F)$: each coordinate measures $X$ against one opposition
$B_i$. The associative series of $X$ is the set of nodes that share a
coordinate with it,
\[
  \{ Y : c_i(Y) = c_i(X) \text{ for some fixed } i \},
\]
grouped by each opposition in turn. The frame $F$ is the memory in which these
groups live; the soft key profile is the relaxed form of membership, valid
also for nodes outside the span.
\end{proposition}

\begin{remark}[In praesentia / in absentia, exactly]
The two readings use disjoint data. The syntagmatic reading uses only what is
present together in the window (the document and the context); the associative
reading uses the frame, whose basis vectors need not be present in the
current window at all. Saussure's distinction is therefore not a nuance of
emphasis but a type distinction between two functions with different domains.
\end{remark}

% ============================================================================
\section{Mechanism, grammar, abstract entities}
\label{sec:mechanism}
% ============================================================================

\subsection{The mechanism of language}
\label{sec:mechanism-lang}

Chapter VI describes the functioning of the system: forms are generated and
interpreted by the interplay of the two kinds of relation
\cite[Part II, ch.~VI]{harris1983}.

\begin{reading}[The mechanism]
\label{read:mechanism}
The mechanism is the composition of the two readings. The dispatcher
\[
  \Disp : \Lat \times \mathrm{State}
  \;\longrightarrow\;
  \mathrm{Dist}(\mathrm{Hosts})
\]
reads the lattice state, which contains both axes: the coverage numbers
$N_D,R_D,\Cov$ (syntagmatic) and the frame profile or hard key (associative).
It then selects which host actualizes the shared langue.

The hosts, however, are not abstract labels. In the multi-LLM scenario from
which the EWM architecture arose, a user works simultaneously with several
LLMs; in the lattice each host appears as a \emph{participant}: a context,
presented not as a separate sub-lattice but as the subset of HLLSets in the
current lattice that the participant's context has produced and refers to.
Two participants understand each other within their \emph{common
context}---the intersection of their context subsets---and may disagree
outside it. The mechanism of language, so read, is the mechanism of a
conversation: each participant speaks from a context, the shared langue is
the common part, and the dispatcher's selection is the turn-taking of
dialogue.
\end{reading}

\begin{definition}[Participants and common context]
\label{def:participants}
Let $H_1, H_2 \subseteq \Lat$ be the context subsets of two participants in
the current fiber. Their \textbf{common context} is the intersection
\[
  H_1 \cap H_2,
\]
the nodes both participants hold together, and their \textbf{disagreement} is
the symmetric difference
\[
  H_1 \xor H_2 \;=\; (H_1 \setminus H_2) \cup (H_2 \setminus H_1).
\]
Mutual understanding is exact on $H_1 \cap H_2$, and only there; outside the
common context the participants may speak past one another.
\end{definition}

\begin{observation}[Understanding is graded]
\label{obs:graded-understanding}
On the common context $H_1 \cap H_2$, the two participants' readings agree
exactly: same nodes, same frame coordinates. Outside it their readings agree
only partially, and the soft key and the spill measure how far apart their
contexts are. Disagreement, like spill, is not noise: it is the boundary of
the shared langue, and the dispatcher's task is to choose the participant
whose context covers the current state best.
\end{observation}

\subsection{Grammar and its subdivisions}
\label{sec:grammar}

Chapter VII draws the consequence that morphology, syntax, and lexicology are
not autonomous domains: they interpenetrate, because all of them are
organizations of the same system of values \cite[Part II, ch.~VII]{harris1983}.

\begin{proposition}[The subdivisions of grammar are one algebra]
\label{prop:grammar}
On one plane $\Plane$, morphology (composition of units), syntax
(combinations of units), and lexicology (values of units) are the same
operations read at different granularities: joins and meets are both word
formation and phrase formation; the ring sum is both morpheme contrast and
phrase contrast; the hard key is both the value of a morpheme and the value
of a construction. The classical boundaries are granularities of the
materialized chain, not types of algebra.
\end{proposition}

\begin{table}[htbp]
\centering
\renewcommand{\arraystretch}{1.15}
\begin{tabular}{@{}>{\raggedright\arraybackslash}p{3.2cm}
                    >{\raggedright\arraybackslash}p{4.4cm}
                    >{\raggedright\arraybackslash}p{7.4cm}@{}}
\toprule
Subdivision & In the \textit{Course} & In HLLSet Algebra \\
\midrule
Morphology & composition and derivation of units & join $\cup$; idempotence
collapses reduplication; ring sum $\xor$ as morpheme contrast \\
Syntax & combinations in the chain & span of the frame as the set of
well-formed constructions; order restored by materialization \\
Lexicology & the system of values & the hard key $c$ as the value; the value
table of Corollary~\ref{cor:lexicon} as the lexicon \\
Boundaries & subdivisions interpenetrate & one Boolean lattice and one GF(2)
ring over one plane: there is no boundary to draw \\
\bottomrule
\end{tabular}
\caption{Grammar re-founded: three subdivisions, one algebra.}
\label{tab:grammar}
\end{table}

\subsection{Abstract entities}
\label{sec:abstract}

Chapter VIII defends the abstract entities of grammar (a genitive, a plural,
a zero ending): they are real because they are positions in the system, not
because they have a sound \cite[Part II, ch.~VIII]{harris1983}.

\begin{reading}[Abstract entities]
\label{read:abstract}
The abstract entities of the algebra are the coordinates themselves: the
basis vectors $B_i$, the hard keys $c$, the generation stamps $\Gen$, and the
frame $F$. They have no tokens and no sound, yet they are the most real
objects of the system, because every concrete sign is a value only through
them. A zero ending is $\varnothing$ as a value: real as a coordinate,
inaudible as a sound.
\end{reading}

\begin{observation}[The frame is an abstract entity]
The frame $F$ never appears as a token in any host's parole; it is recomputed
from observations and exists only as the coordinate system of the fiber.
Saussure's chapter VIII is the recognition, seventy years before the
datatype, that a linguistic system contains objects that have reality but no
substance. The generation-stamped frame is exactly such an object.
\end{observation}

% ============================================================================
\section{The dictionary of Part Two}
\label{sec:dictionary}
% ============================================================================

\begin{table}[htbp]
\centering
\renewcommand{\arraystretch}{1.15}
\begin{tabular}{@{}>{\raggedright\arraybackslash}p{3.6cm}
                    >{\raggedright\arraybackslash}p{4.4cm}
                    >{\raggedright\arraybackslash}p{6.8cm}@{}}
\toprule
Saussure (Part Two) & In the \textit{Course} & In HLLSet Algebra \\
\midrule
concrete entity & segment uniting concept and sound pattern & the sign pair
$(\Ing,\Mat)$; delimitation by the coverage round-trip \\
delimitation of units & cutting the continuous chain & coverage
$\Cov(\Ctx)=1$; spill measures the failure \\
identity & the $8{:}45$ express & content addressability:
Proposition~\ref{prop:identity} \\
reality & position in the system & position in the relational structure of the
frame; hard key as its compressed form; soft key plus spill outside the span \\
value (conceptual) & opposition within a system & relational position among
nodes: Reading~\ref{read:conceptual}; hard key $c\in\mathrm{GF}(2)^k$ as its
frame compression: Theorem~\ref{thm:complete}; similar concepts need two
bijections plus matching co-occurrence: Observation~\ref{obs:similar-concepts} \\
value (material) & phonic differences & symmetric difference $\xor$:
Proposition~\ref{prop:material-difference} \\
value (whole sign) & only differences, no positive terms &
Theorem~\ref{thm:differences} \\
syntagmatic relation & in praesentia, in the chain & $\Synt(X;D)=(N,R,\Cov)$
along the materialized chain \\
associative relation & in absentia, in memory & $\Ass(X;F)=(\BSS_{B_i}(X))$;
hard key on the span \\
mechanism of language & interplay of the two relations & dispatcher composing
both readings: Reading~\ref{read:mechanism}; participants as context subsets,
understanding on their intersection: Definition~\ref{def:participants} \\
grammar's subdivisions & morphology, syntax, lexicology & one algebra, three
granularities: Proposition~\ref{prop:grammar} \\
abstract entities & genitive, plural, zero ending & basis vectors, hard keys,
generation stamps, the frame itself \\
\bottomrule
\end{tabular}
\caption{The complete dictionary of Part Two of the \textit{Course} into
HLLSet Algebra.}
\label{tab:dictionary}
\end{table}

% ============================================================================
\FloatBarrier
\section{What Part II adds to Part I}
\label{sec:adds}
% ============================================================================

Part I promised value and did not deliver it; Part II delivers it and thereby
closes the synchronic half of the exegesis. Four additions matter.

\begin{enumerate}[leftmargin=*,nosep]
  \item \textbf{Value is now a datatype, not a promise.} Value proper is a
    relational structure $(N,R)$ with co-occurrence profiles
    (Definition~\ref{def:relational}); its exact compression on the span is
    the hard key, a vector over GF(2); the relaxed form is the soft key, a
    vector over $[0,1]$; the spill is a cardinal. All of these can be stored,
    compared, and generation-stamped.
  \item \textbf{Identity and value are proven to coincide.} Within a fiber,
    Theorem~\ref{thm:complete} turns Saussure's chapter III claim (identity is
    position in the system) into an equivalence, not an aphorism.
  \item \textbf{``Differences without positive terms'' is a theorem.}
    Theorem~\ref{thm:differences} is the sentence of the \textit{Course} made
    literal: a node is its difference pattern, and the only positive terms
    (the basis vectors) are arbitrary by Part I's coordinate freedom.
  \item \textbf{The two axes of Part One and the two relations of Part Two
    are separated.} The AB/CD axes of Part One are temporal (fiber vs
    trajectory); the syntagmatic/associative relations of Part Two are both
    \emph{inside} one fiber (chain vs frame). Confusing the two distinctions
    is a known hazard of Saussure commentary; the algebra makes the confusion
    impossible because the two pairs have different types.
\end{enumerate}

% ============================================================================
\section{Conclusion and outlook to Part III}
\label{sec:conclusion}
% ============================================================================

Part Two of the \textit{Course} has been re-founded as the theory of
relational value, compressed without loss on the span by the hard key.
Concrete entities are sign pairs delimited by coverage; identity is
content addressability; reality and value are position in the frame; the two
relations are the two readings, syntagmatic and associative; the mechanism is
the dispatcher; grammar is one algebra over one plane; and the abstract
entities are the coordinates themselves. Saussure's synchronic linguistics,
read in this direction, is no longer a philosophical description of a system
of values: it is the specification of a computable one, and the reference
implementation satisfies it.

What Part Two deliberately excludes is time. The frame was held fixed; the
spill was treated as a boundary, not as a cause. Part III will release that
exclusion: diachronic linguistics, in which phonetic change and analogy become
basis pivots and generation drift, the diachronic law (imperative but not
general) becomes the recorded transition, and Saussure's distinction between
the two perspectives is tested against a system in which the frame actually
moves. The residual handed to Part III is the one Part II named but did not
study: the pivot itself.

\section*{Acknowledgments}
This paper was created in collaboration with, and with the assistance of, an
AI system---DeepSeek (deepseek-v4-pro). The hard-key reading of value, the
completeness theorem, and several of the formulations emerged in dialogue
with the assistant, whose drafts the authors revised and take full
responsibility for in their final form.

\begin{thebibliography}{99}

\bibitem{saussure1916} F.~de Saussure,
\textit{Cours de linguistique g\'en\'erale}, Payot, Paris, 1916.
(Edited posthumously by C.~Bally and A.~Sechehaye, with A.~Riedlinger.)

\bibitem{harris1983} F.~de Saussure,
\textit{Course in General Linguistics}, translated and annotated by
Roy Harris, Duckworth, London, 1983.

\bibitem{baskin1959} F.~de Saussure,
\textit{Course in General Linguistics}, translated by Wade Baskin,
Philosophical Library, New York, 1959.

\bibitem{saussure1977} F.~de Sossyur,
\textit{Trudy po yazykoznaniyu} [Works in Linguistics], edited by
A.~A.~Kholodovich, translation of the \textit{Cours} by A.~M.~Sukhotin,
Progress, Moscow, 1977.
(Authors' working copy; used for verification of the readings.)

\bibitem{part1} A.~Mylnikov, A.~Mylnikov Jr., and D.~Kumar,
\textit{Saussure in HLLSet Algebra, Part I: General Principles},
SAUSSURE PUB series, 2026.

\bibitem{ewm2026} A.~Mylnikov, A.~Mylnikov Jr., and D.~Kumar,
\textit{Emerging World Models: A Structural Framework for Self-Organizing
World Representations}, EWM PUB series, 2026.

\bibitem{companion2026} A.~Mylnikov, A.~Mylnikov Jr., and D.~Kumar,
\textit{The EWM Lattice and Its Worlds}, EWM PUB series, 2026.

\bibitem{state2026} A.~Mylnikov, A.~Mylnikov Jr., and D.~Kumar,
\textit{Transformations and Trajectories of a Content-Addressed State
Machine}, EWM PUB series, 2026.

\bibitem{attention2026} A.~Mylnikov, A.~Mylnikov Jr., and D.~Kumar,
\textit{Attention on Lattice Traversal Is All You Need}, EWM PUB series, 2026.

\bibitem{langue2026} A.~Mylnikov, A.~Mylnikov Jr., and D.~Kumar,
\textit{Langue as a Lattice: Closing the Circle from Statistical to
Structural Linguistics through Content-Addressed Algebras}, EWM PUB series,
2026.

\bibitem{birkhoff1940} G.~Birkhoff,
\textit{Lattice Theory}, American Mathematical Society, 1940.

\bibitem{maclane1971} S.~Mac Lane,
\textit{Categories for the Working Mathematician}, Springer, 1971.

\bibitem{hjelmslev1961} L.~Hjelmslev,
\textit{Prolegomena to a Theory of Language}, University of Wisconsin Press,
1961.

\bibitem{harris1954} Z.~S.~Harris,
\textit{Distributional Structure}, Word, 10(2--3):146--162, 1954.

\bibitem{firth1957} J.~R.~Firth,
\textit{A Synopsis of Linguistic Theory, 1930--1955}, Studies in Linguistic
Analysis, 1957.

\bibitem{vaswani2017} A.~Vaswani, N.~Shazeer, N.~Parmar, J.~Uszkoreit,
L.~Jones, A.~N.~Gomez, \L.~Kaiser, and I.~Polosukhin,
\textit{Attention Is All You Need}, NeurIPS, 2017.

\end{thebibliography}

\end{document}
